\subsection{Modular symmetry and non-Abelian discrete flavor symmetries in string compactification --- arXiv:1804.06644}

\textbf{Authors:} Tatsuo Kobayashi, Satoshi Nagamoto, Shintaro Takada, Shio Tamba, Takuya H. Tatsuishi

\paragraph{主な主張}
磁化した $T^2$ 上の D-brane 模型では、磁束 $M=2$ の場合に非可換フレーバー対称性 $D_4$ がモジュラー対称性の部分群として現れ、heterotic $T^2/Z_4$ orbifold でも同じモジュラー対称性と $D_4$ の包含関係が実現することを示した \cite{Kobayashi:2018rad}。

\paragraph{新規性}
従来の非可換離散フレーバー対称性とモジュラー対称性を独立な構造として扱うのではなく、具体的な弦コンパクト化において前者が後者の部分群として統一的に理解できる例を明示した点にある。

\paragraph{位置付け}
磁化 D-brane・heterotic orbifold におけるフレーバー対称性の幾何学的起源とモジュラー対称性を結び付ける top-down 研究であり、後のモジュラーフレーバー模型を弦理論側から理解するための基礎的な位置を占める。
